\section{Reconstrucción incidencial y covariancia}
\label{nuclear:5}
\subsection{La isometría completa y su dominio}\label{nuc05:isometria}

Sea \(H=\mathbb R^{12}\), con producto euclídeo; \(P_0=I-\mathbf1\mathbf1^*/12\), \(\mu(k)=\mathbf1^*k/12\), y \(\mathcal I:H\to\mathbb R^{132}\) la matriz de incidencia de \cref{exc:entrelazador,exc:isometria}. Escribimos \(E_{\rm inc}=\mathcal I(\mathbf1^\perp)\) y

\[
Y=\mathbb R\oplus E_{\rm inc},\qquad
\langle(\mu,y),(\nu,z)\rangle_Y=12\mu\nu+\langle y,z\rangle.
\]

La isometría completa construida en la sección~\ref{nuclear:2} es

\[
F:H\longrightarrow Y,\qquad Fk=(\mu(k),\mathcal I P_0k/6),
\]

\[
F^{-1}(\mu,y)=\mu\mathbf1+P_0\mathcal I^*y/6.
\]

El Gram \(\mathcal I^*\mathcal I=36I+30\mathbf1\mathbf1^*\) prueba ambas composiciones inversas y \(F^*F=I_H\), \(FF^*=I_Y\), con el producto ponderado indicado. Por ello \(F\) es ortogonal entre \(H\) e \(Y\). Puede complejificarse sin cambiar las pruebas.

El codominio es \(Y\), de dimensión doce, no todo \(\mathbb R\oplus\mathbb R^{132}\). Extender un operador desde \(Y\) a su complemento en ese ambiente requiere una elección adicional; no forma parte de la dinámica transportada de manera única. En coordenadas firmadas normalizadas se usa \(FJ_H\), con \(J_H=H_{12}/2\). La inversión integral sigue siendo \(H_{12}/4\) en su subred, no un normalizador métrico alternativo.

\subsection{Teorema de levantamiento covariante en la imagen}\label{nuc05:levantamiento}

\begin{theorem}
Para un operador acotado \(D\) sobre \(H\), existe un único operador sobre \(Y\) que satisface \(D_YF=FD\). Es

\[
D_Y=FDF^{-1}.
\]

La correspondencia conserva sumas, productos, adjuntos, norma, positividad e identidad. En particular, lleva unitarios a unitarios y observables autoadjuntos a observables autoadjuntos. Para un unitario \(C\) y un observable \(A\),

\[
C^*AC=A\quad\Longleftrightarrow\quad C_Y^*A_YC_Y=A_Y.
\]

Si \(R\) es un cambio unitario de carta, \(R_Y=FRF^{-1}\), entonces

\[
F(RDR^{-1})F^{-1}=R_YD_YR_Y^{-1}.
\]
\end{theorem}

\begin{proof}
La primera fórmula se obtiene multiplicando el entrelazamiento por \(F^{-1}\). La unicidad se sigue de la sobreyectividad hacia \(Y\). Las identidades de productos y adjuntos resultan de \(F^{-1}=F^*\), y la igualdad de normas de la isometría sobreyectiva. La última igualdad es una composición asociativa de esas aplicaciones. No interviene ninguna hipótesis de conmutación entre \(R\) y \(D\).
\end{proof}

Para \(g\in\operatorname{Aut}(\mathcal H)\), la proposición~\ref{exc:isometria} proporciona

\[
F\rho_{12}(g)=\bigl(1\oplus\rho_{\rm hex}(g)\bigr)F.
\]

Así, el cambio de carta posee una acción combinatoria concreta sobre las hexadas. En cambio, para un unitario general \(D\), \(D_Y\) es un operador de la representación incidencial, no necesariamente una permutación del diseño. Preservar el producto interno de esa representación no convierte a todos sus unitarios en automorfismos combinatorios ni en elementos del Monstruo.

\subsubsection{Covariancia bajo un cambio de carta no conmutante}

Sea \(C\) el avance \(Ce_p=e_{p+1}\) cíclico en las doce posiciones, es decir, el sentido inverso de la convención \(S\) de \cref{sec:k-registro}. Es una operación de la carta de ventanas, no la holonomía temporal completa. Sea \(R\) la estrella ya construida

\[
R=(1\ 3)(2\ 11)(5\ 7)(8\ 12).
\]

Entonces \(RCe_1=e_{11}\), mientras \(CRe_1=e_4\). No conmutan. Sin embargo, para \(C'=RCR^{-1}\), se cumplen exactamente \(T(C')R=RT(C)\) y \(\eta(C')R=R^{\oplus9}\eta(C)\), y lo mismo en \(Y\) mediante \(F\). La covariancia de carta sí funciona en un ejemplo donde la conmutación en carta fija falla.

\subsection{Recuperación explícita desde el registro incidencial de complementos}\label{nuc05:recuperacion}

Sean \(C\) unitario, \(T=(8I+C)/9\), \(\eta=(I-JJ^*)S_CJ:H\to H^9\), y \(P\) la proyección sobre \(\operatorname{Fix}(C)\), como en la sección~\ref{nuclear:2}. Escribimos \(F_9=F^{\oplus9}\). La versión incidencial del registro finito es

\[
\mathcal W_Nk=\left(FT^Nk,F_9\eta k,F_9\eta Tk,\ldots,F_9\eta T^{N-1}k\right).
\]

\begin{theorem}
Esta aplicación es isométrica y posee la inversa izquierda explícita

\[
\mathcal W_N^*(a,b_0,\ldots,b_{N-1})
=T^{*N}F^{-1}a+
\sum_{j=0}^{N-1}T^{*j}\eta^*F_9^{-1}b_j.
\]

En particular, aplicar esta fórmula a \(\mathcal W_Nk\) recupera exactamente \(k\). El registro infinito

\[
\mathcal W_\infty k=\left(FPk,(F_9\eta T^jk)_{j\ge0}\right)
\]

también es isométrico, y su inversa izquierda es

\[
k=PF^{-1}a+\sum_{j\ge0}T^{*j}\eta^*F_9^{-1}b_j
\quad\text{para }(a,(b_j))=\mathcal W_\infty k.
\]

La serie converge en norma. Para datos arbitrarios cuadrado-sumables la misma fórmula define el adjunto acotado, aunque no todo dato del codominio pertenece a la imagen de \(\mathcal W_\infty\).
\end{theorem}

\begin{proof}
Aplicar \(F\) y \(F_9\) a cada componente de \(V_N\) o \(V_\infty\) conserva el producto interno. Tomar el adjunto de esa composición da las fórmulas. La identidad telescópica de la sección~\ref{nuclear:2} prueba \(\mathcal W_N^*\mathcal W_N=I\), y su límite fuerte prueba \(\mathcal W_\infty^*\mathcal W_\infty=I\). La convergencia de la serie procede de aplicar un operador acotado a las truncaciones de una sucesión en suma directa hilbertiana.
\end{proof}

Esta reconstrucción no conserva \(k\) como una primera coordenada intacta: el primer componente es el terminal comprimido o su sector fijo; el resto procede de los complementos sucesivos. Cuando \(P=0\), los complementos incidenciales por sí solos recuperan \(k\). La contribución recuperada a partir de los primeros \(N\) bloques es

\[
k_N=\sum_{j<N}T^{*j}\eta^*F_9^{-1}b_j
=\bigl(I-T^{*N}T^N\bigr)k.
\]

El residuo es \(T^{*N}T^Nk\), que tiende a cero si \(P=0\). Con un sector fijo no nulo debe conservarse \(FPk\); la memoria de complementos sola no recupera ese sector. Las tasas de reconstrucción y la dinámica sobre el índice de bloques están desarrolladas en la sección~\ref{nuclear:3}.

Si \(A\) es autoadjunto y \([A,C]=0\), la forma del observable se conserva también:

\[
\langle k,Ak\rangle
=\langle FPk,A_YFPk\rangle_Y
+\sum_{j\ge0}\langle F_9\eta T^jk,(I_9\otimes A_Y)F_9\eta T^jk\rangle.
\]

La serie escalar converge absolutamente porque \(A\) es acotado y la suma de las normas cuadradas de los componentes es finita. La hipótesis expresa conservación de \(A\) por \(C\); no exige conmutación de \(C\) con los cambios de carta excepcionales.

\subsection{Levantamiento natural del refinamiento con marcos locales}\label{nuc05:refinamiento}

La construcción de \cref{iv:cont:historias,iv:cont:medida} aporta conjuntos de estados cilíndricos enriquecidos \(Z_n\), truncamientos y medidas proyectivas de soporte completo. Para un descendiente \(z'\) de \(z\),

\[
w_n(z'|z)=\sqrt{\mu_{n+1}(z')/\mu_n(z)},\qquad
\sum_{\tau z'=z}w_n(z'|z)^2=1.
\]

La etiqueta \(z'\) conserva registro ordenado, orientación, acarreo, frontera, memoria, ruta y terminal. El refinamiento de un estado no identifica descendientes distintos.

Sean \(q(z)\) el registro entero de memoria, \(C\) un unitario sobre \(H\), y \(R_z\) cambios ortogonales de marco de la fibra. Para marcos incidenciales concretos puede usarse \(R_z=\rho_{12}(g_z)\), con \(g_z\in\operatorname{Aut}(\mathcal H)\), transportando sus marcas. Definimos

\[
G_z=R_zC^{q(z)},\qquad
U_{z',z}=G_{z'}G_z^*
=R_{z'}C^{q(z')-q(z)}R_z^*.
\]

Esta construcción prolonga la equivalencia de calibre explícita de \cref{exc:doble-lectura}. No requiere \([R_z,C]=0\), y deja visible el orden de las operaciones.

Sobre \(\mathscr H_n=\ell^2(Z_n)\otimes H\),

\[
\mathscr U_n(e_z\otimes v)
=\sum_{\tau z'=z}w_n(z'|z)e_{z'}\otimes U_{z',z}v.
\]

Definimos \(\mathscr F_n=I_{\ell^2(Z_n)}\otimes F\), con codominio \(\mathscr Y_n=\ell^2(Z_n)\otimes Y\).

\begin{theorem}
El refinamiento es isométrico y satisface

\[
\mathscr U^Y_n\mathscr F_n=\mathscr F_{n+1}\mathscr U_n,
\quad
\mathscr U^Y_n=\mathscr F_{n+1}\mathscr U_n\mathscr F_n^{-1}.
\]

La fórmula de cada término es la anterior con \(U^Y_{z',z}=FU_{z',z}F^{-1}\). Los sistemas de varias profundidades conservan composición y admiten una isometría sobreyectiva entre sus límites inductivos. Se preservan exactamente todos los coeficientes de fibra con sus etiquetas \(z\).
\end{theorem}

\begin{proof}
Los descendientes de padres distintos son ortogonales; cada \(U_{z',z}\) es unitario; y los pesos tienen suma cuadrática uno. Esto prueba la isometría. En una trayectoria, los factores intermedios \(G_z^*G_z\) se cancelan y los cocientes de medidas telescopan. La composición no depende de la partición de la trayectoria. El cuadrado con \(F\) se verifica sobre \(e_z\otimes v\), término a término. Las aplicaciones y sus inversas son isometrías compatibles, por lo que se extienden inversamente a las completaciones inductivas.
\end{proof}

\subsubsection{Observables y dinámica compatibles}

Si \(A_0\) es un observable acotado de la fibra, sea

\[
(\mathscr A_n\psi)_z=G_zA_0G_z^*\psi_z.
\]

Entonces \(\mathscr A_{n+1}\mathscr U_n=\mathscr U_n\mathscr A_n\), de donde \(\mathscr U_n^*\mathscr A_{n+1}\mathscr U_n=\mathscr A_n\). Se obtiene sustituyendo \(U_{z',z}=G_{z'}G_z^*\). Esta es conservación de un campo de observables transportado por refinamiento; no afirma que \(A_0\) conmute con \(C\) en una carta fija.

La dinámica unitaria por nivel

\[
(\mathscr C_n\psi)_z=G_zCG_z^*\psi_z
\]

satisface \(\mathscr C_{n+1}\mathscr U_n=\mathscr U_n\mathscr C_n\). También entrelazan sus adjuntos: como ambas \(\mathscr C\) son unitarias, la identidad se multiplica por sus inversas. Por tanto entrelazan \(T\), sus complementos nonádicos y sus proyecciones fijas. El registro finito y el infinito de la sección~\ref{nuc05:recuperacion} son naturales con respecto a este refinamiento. Al aplicar \(\mathscr F_n\), esos mismos cuadrados valen en las fibras incidenciales. Para que \(\mathscr A_n\) sea conservado además por esa dinámica en un nivel fijo, la condición precisa es \([A_0,C]=0\); la covariancia entre marcos sigue sin requerir \([R_z,C]=0\).

\subsubsection{Especialización de memoria bilateral con marcos variables}

En la órbita reversible de \cref{nuclear:3} puede representarse una publicación de coordenadas por \(\psi=(\psi_m)_{m\in\mathbb Z}\in\ell^2(\mathbb Z;H)\). Su avance es \((C_0\psi)_m=\psi_{m-1}\). \(F\) actúa punto a punto. Para marcos \(R_m=\rho_{12}(g_m)\), el cambio unitario \((\mathscr R\psi)_m=R_m\psi_m\) produce

\[
(C'\psi)_m=R_mR_{m-1}^*\psi_{m-1}.
\]

El operador incidencial tiene la misma fórmula con \(1\oplus\rho_{\rm hex}(g_mg_{m-1}^{-1})\). No se imponen marcos constantes ni conmutación con el desplazamiento. Como \(C'\) es conjugado a \(C_0\), no tiene vectores fijos no nulos en este espacio cuadrado-sumable. La sección~\ref{nuc05:recuperacion} recupera la publicación lineal completa desde sus complementos incidenciales. Se trata de la órbita y de las coordenadas representadas: no de una identificación de todas las historias TPK con \(\mathbb Z\).

\subsection{Extensión directa a las doce fibras \texorpdfstring{\(A_2\)}{A2}}\label{nuc05:a2}

La construcción excepcional conserva doce fibras \(A_2\) y el operador \(g_c\) de orden tres desarrollado en la sección~\ref{nuc04:c3}. Sea \(V_{A_2}=A_2\otimes_{\mathbb Z}\mathbb R\) su realización euclídea bidimensional. La misma isometría posicional puede tensorizarse sobre \(\mathbb R\) con ese espacio:

\[
F_{A_2}=F\otimes I_{V_{A_2}}:
\mathbb R^{12}\otimes_{\mathbb R}V_{A_2}
\longrightarrow Y\otimes_{\mathbb R}V_{A_2}.
\]

Su primera componente es la media vectorial y su segunda son las sumas incidenciales centradas, ahora con valores en \(V_{A_2}\). El Gram de incidencia, tensorizado con la identidad, prueba de nuevo la isometría completa y la inversa. Por ello \(g_c^Y=F_{A_2}g_cF_{A_2}^{-1}\) es ortogonal, de orden tres y fijo-libre. La composición nonádica de la sección~\ref{nuclear:2} puede publicarse y recuperarse en esta pantalla con valores en \(V_{A_2}\), mediante las secciones~\ref{nuc05:levantamiento} y~\ref{nuc05:recuperacion}. No se confunden los doce escalares de \(K\) con las veinticuatro coordenadas de las fibras \(A_2\), ni se convierte la acción en una permutación de hexadas cuando la elevación también actúa dentro de las fibras.

\subsection{Alcance exacto respecto de las historias}

Sea \(k:X\to H\) el lector del registro sobre un dominio de historias, y sea \(\ell=F\circ k\). Entonces

\[
\ell(x)=\ell(y)\quad\Longleftrightarrow\quad k(x)=k(y).
\]

La pantalla completa añade cero pérdida respecto de \(K\), pero tampoco cambia las fibras de ese lector. Un mapa \(d:X\to X\) desciende a una dinámica sobre los valores de \(K\) si y sólo si

\[
k(x)=k(y)\Longrightarrow k(d x)=k(d y).
\]

La misma condición es necesaria y suficiente para descender sobre la imagen incidencial. La prueba consiste en definir \(D(k(x))=k(dx)\); la condición equivale a que la definición sea independiente del representante. Si ese \(D\) es lineal y unitario en la realización declarada, se aplica la sección~\ref{nuc05:levantamiento}. Si es no lineal, \(F\) sigue conjugándolo como aplicación, pero no lo convierte en un unitario.

El criterio distingue el descenso a \(K\) de la dinámica enriquecida. La sección~\ref{subsec:k-terminal} establece que el lector utiliza subruta, orientación, acarreo e incidencias; la sección~\ref{exc:limite-inverso} distingue recuperación de palabra visible, soporte y memoria profunda. En \(\ell^2(Z_n)\otimes H\), la sección~\ref{nuc05:refinamiento} conserva las etiquetas enriquecidas y transporta la representación de la fibra. La equivalencia actúa sobre el dominio de historias previamente generado.

Para un corte ya construido, la naturalidad del registro en \cref{prop:k-extension} fija \(K_N=K_{108}\circ\tau_{N,108}\): \(F\) conserva exactamente esa naturalidad. Una sucesión de nuevos eventos puede actualizar \(K\); la composición aditiva o afín de \cref{sec:k-registro} conserva la independencia de la partición, no la constancia del valor al añadir eventos.
